\begin{abstract}
We present NeuroGenUI, a multimodal framework that generates editable React user interfaces from a hand-drawn sketch and spoken instructions, and refines them with further voice commands. Instead of a single image-to-code call, the system separates perception, interpretation, fusion, generation and validation around an explicit JSON UI Intent Representation: OpenCV proposals are labelled by a vision model, speech is transcribed by Whisper and converted into typed edit operations, a deterministic geometry-aware compiler emits React code, and a headless-browser validator drives a targeted repair loop. We evaluate every stage on \CountSketches{} labelled synthetic sketches and \CountCommands{} synthetic spoken commands, against a direct sketch-to-UI baseline using the same model and measurement code. On rendered UIs, the pipeline detects the sketched components with F1 \ComponentDetectionFODomPipeline{}\% versus \ComponentDetectionFODomBaseline{}\%, passes \FunctionalCorrectnessPipeline{}\% of specified interactions versus \FunctionalCorrectnessBaseline{}\%, and applies refinements in \ModificationTimeMsPipeline{}\,s versus \ModificationTimeMsBaseline{}\,s while preserving \PreservationRatePipeline{}\% of unaffected components; requirement satisfaction and refinement success are at parity with the baseline. Fault injection shows the validator detects \FaultDetectionRatePipeline{}\% of injected faults and the repair loop restores \FaultRepairRatePipeline{}\%. An initial evaluation exposed defects that were fixed on the same test set, so results are development-set figures on synthetic data. Code, dataset and evaluation scripts are provided.
\end{abstract}
